% sections/05_normalization.tex  –  Normalization (Stage 2)
% \input{} into main.tex; do NOT wrap in \documentclass or \begin{document}

\section{Normalization}
\label{sec:norm}

Normalization is \textbf{Stage~2} of the data-conditioning pipeline
introduced in \Cref{sec:preproc}.  Its purpose is to rescale feature
values so that learning algorithms converge faster and are less
sensitive to the raw physical units of the input channels.  Every
normalization method discussed here obeys a single inviolable rule:
\emph{statistics are computed on the training partition only and applied
identically to validation and test partitions.}

%% ------------------------------------------------------------------
\subsection{Anti-Leakage Principle}
\label{sec:norm:leakage}

\begin{tcolorbox}[
  colback=red!5!white,
  colframe=red!75!black,
  title={\textbf{Warning — Data-Leakage Risk}},
  fonttitle=\bfseries
]
Fitting normalization statistics (mean, standard deviation, minimum,
maximum, median, or interquartile range) on the \emph{entire} dataset —
including validation and test partitions — constitutes \textbf{data
leakage}.  The resulting performance estimates are optimistically biased
because the model indirectly observes distributional information from
held-out samples during training.  All normalization parameters must be
computed exclusively on the training fold and then frozen before
transforming any held-out data.
\end{tcolorbox}

In a cross-validation setting (\Cref{sec:overfit}), this means
normalization statistics are \emph{re-computed} for every fold: the
inner training set defines the transform, and the inner validation set
(and outer test set, in nested CV) receives the same frozen transform.

%% ------------------------------------------------------------------
\subsection{Per-Feature Z-Score Normalization}
\label{sec:norm:zscore}

Z-score (standard-score) normalization centres each feature at zero mean
and unit variance:
\begin{equation}
  x'_j = \frac{x_j - \mu_j^{\text{train}}}
               {\sigma_j^{\text{train}}},
  \label{eq:zscore}
\end{equation}
where $j$ indexes the feature (i.e., the frequency bin or pixel
position), and $\mu_j^{\text{train}}$, $\sigma_j^{\text{train}}$ are
the mean and standard deviation computed over the training partition.

\paragraph{When to prefer z-score.}
Z-score normalization is appropriate when the feature distributions are
approximately Gaussian or at least unimodal and symmetric.  It is the
default choice for algorithms that assume or benefit from standardised
inputs, including Logistic Regression~\cite{hosmer2013logistic},
Linear SVM~\cite{vapnik1995svm,cortes1995svm}, and Multilayer
Perceptrons~\cite{haykin2009nn,goodfellow2016deep}.

\paragraph{Sensitivity to outliers.}
Because $\mu$ and $\sigma$ are influenced by extreme values, z-score
normalization can compress the bulk of the distribution when heavy-tailed
outliers are present.  If the outlier inspection in \Cref{sec:preproc}
reveals residual anomalies, robust scaling (\Cref{sec:norm:robust}) is
preferred.

%% ------------------------------------------------------------------
\subsection{Min-Max Scaling}
\label{sec:norm:minmax}

Min-max scaling maps each feature linearly into the interval $[0,\,1]$:
\begin{equation}
  x'_j = \frac{x_j - x_{j,\min}^{\text{train}}}
               {x_{j,\max}^{\text{train}} - x_{j,\min}^{\text{train}}}.
  \label{eq:minmax}
\end{equation}

\paragraph{When to prefer min-max.}
Min-max scaling is useful when a bounded output range is desirable, for
example as input to layers with sigmoid or softmax activations, or when
converting traces to greyscale pixel values in $[0,\,255]$ for
Route~B image representations (\Cref{sec:repr}).  It preserves zero
entries and the shape of the original distribution, but it is
\emph{highly} sensitive to outliers: a single extreme value can compress
all other values into a narrow sub-interval.

\paragraph{Out-of-range values at test time.}
Because the min and max are determined from training data, test-time
values may fall outside $[0,\,1]$.  The recommended policy is to
\emph{clip} to $[0,\,1]$ and log any clipped samples rather than
silently extrapolating.

%% ------------------------------------------------------------------
\subsection{Robust Scaling}
\label{sec:norm:robust}

When outlier contamination is suspected — or confirmed by the $4\sigma$
screening in \Cref{sec:preproc} — robust scaling replaces mean and
standard deviation with the median ($\tilde{x}$) and interquartile range
(IQR):
\begin{equation}
  x'_j = \frac{x_j - \tilde{x}_j^{\text{train}}}
               {\text{IQR}_j^{\text{train}}},
  \qquad
  \text{IQR}_j = Q_{75,j} - Q_{25,j}.
  \label{eq:robust}
\end{equation}

Robust scaling does not bound the output range, but it ensures that the
central $50\,\%$ of the data occupies a unit-width interval regardless
of tail behaviour.  It is compatible with all algorithms considered in
this report and is recommended as a default for exploratory runs before
the outlier landscape is fully characterised.

%% ------------------------------------------------------------------
\subsection{Channel-Wise Normalization for Route~A}
\label{sec:norm:channelwise}

Route~A (1D spectral) represents each sample as a dual-channel vector:
one channel for $|S_{11}|$~(dB) and one for $\angle S_{11}$~(deg).
The two channels have fundamentally different scales — magnitude values
typically range from \SI{0}{\decibel} to \SI{-40}{\decibel}, while
phase values span $\pm\SI{180}{\degree}$ — and must therefore be
normalised \emph{independently}.

Concretely, if a sample is represented as
$\mathbf{x} = [\mathbf{m};\,\mathbf{p}] \in \mathbb{R}^{2 \times 4001}$,
the normalization statistics $(\mu,\sigma)$ or $(x_{\min},x_{\max})$ are
computed separately for the magnitude slice $\mathbf{m}$ and the phase
slice $\mathbf{p}$ across the training partition.  This is the
one-dimensional analogue of channel-wise image normalization and is
standard practice in multi-variate time-series
classification~\cite{goodfellow2016deep}.

%% ------------------------------------------------------------------
\subsection{Image Normalization for Route~B}
\label{sec:norm:image}

Route~B converts each $S_{11}$ trace into a 2D image-like
representation (e.g., a magnitude--frequency heatmap or a
recurrence-plot image).  Normalization follows the conventions
established in the computer-vision literature:

\begin{itemize}
  \item \textbf{Single-channel images:} compute pixel-level mean
        $\mu_{\text{px}}$ and standard deviation $\sigma_{\text{px}}$
        across all training images, then apply
        \Cref{eq:zscore} pixel-wise.
  \item \textbf{Dual-channel images} (magnitude and phase as separate
        channels): compute $(\mu,\sigma)$ per channel across the
        training set, paralleling standard RGB-channel normalization.
\end{itemize}

Batch normalization~\cite{ioffe2015batchnorm} provides an additional,
learnable normalization layer inside deep networks.  It normalises
activations within each mini-batch during training and uses running
estimates at inference time.  While batch normalization is architecture-
internal rather than a preprocessing step, it interacts with the input
normalization: applying both z-score input normalization and batch
normalization is common and generally beneficial for training stability.

%% ------------------------------------------------------------------
\subsection{Comparison of Normalization Strategies}
\label{sec:norm:comparison}

\Cref{tab:norm_comparison} summarises the key properties of each
normalization method.

\begin{table}[ht]
\centering
\caption{Comparison of normalization strategies.}
\label{tab:norm_comparison}
\small
\renewcommand{\arraystretch}{1.15}
\begin{tabular}{p{2.4cm} p{2.8cm} p{2.0cm} p{2.0cm} p{2.3cm}}
\toprule
\textbf{Method} &
\textbf{Assumptions} &
\textbf{Outlier sensitivity} &
\textbf{Output range} &
\textbf{Route compat.} \\
\midrule
Z-score &
  Approx.\ Gaussian features &
  High &
  Unbounded &
  A, B \\
Min-max &
  Bounded, no extreme outliers &
  Very high &
  $[0,\,1]$ &
  A, B \\
Robust (median/IQR) &
  Arbitrary; tolerates tails &
  Low &
  Unbounded &
  A, B \\
Batch norm\textsuperscript{$\dagger$} &
  Internal to deep nets &
  Moderate &
  Learned &
  A (1D CNN, ResNet-1D), B (2D CNN) \\
\bottomrule
\multicolumn{5}{l}{\footnotesize $\dagger$\;Batch normalization is an
  architecture component, not a preprocessing step~\cite{ioffe2015batchnorm}.}
\end{tabular}
\end{table}

%% ------------------------------------------------------------------
\subsection{Practical Recommendations}
\label{sec:norm:recommendations}

\begin{enumerate}
  \item \textbf{Default:} begin with per-feature z-score normalization
        (\Cref{eq:zscore}) for all Route~A experiments and pixel-level
        z-score for Route~B.  It is well-understood, widely supported by
        scikit-learn~\cite{sklearn_docs} and TensorFlow, and compatible
        with every algorithm in \Cref{sec:algorithms}.

  \item \textbf{If outliers persist} after the screening in
        \Cref{sec:preproc}, switch to robust scaling
        (\Cref{eq:robust}).  Compare validation performance against
        z-score to confirm the benefit.

  \item \textbf{If bounded inputs are required} (e.g., for sigmoid
        output layers or greyscale image encoding), use min-max scaling
        (\Cref{eq:minmax}) with post-hoc clipping at test time.

  \item \textbf{Always verify:} after normalization, inspect the
        per-feature distributions (histograms or Q--Q plots) on a
        held-out fold to confirm that no pathological compression or
        inflation has occurred.

  \item \textbf{Record the transform:} serialise the fitted scaler
        object (e.g., via \texttt{joblib} or ONNX~\cite{onnx_runtime})
        so that identical normalization can be applied at inference time
        on the edge device.
\end{enumerate}
